\setcounter{section}{4}\setcounter{subsection}{0}\setcounter{prop}{0}
% Original author-published.tex lines 1391--1410
\subsection{Gromov hyperbolic groups and cone types}\label{cont}
Let $\G=\langle S\rangle$ be a finitely generated non-elementary Gromov hyperbolic group, and recall from Section~\ref{cocycles} that we denote by $\bord^2\G$ the set of distinct pairs in its Gromov boundary~$\bord\G$.



\begin{defi}\label{conicalsequence}
A divergent sequence $\{\g_n\}_{n\in\N}\subset\G$ converges to a point $x\in\bord\G$ \emph{conically} if for every $y\in\bord\G-\{x\}$ the sequence $(\g_n^{-1}y,\g_n^{-1}x)$ remains on a~compact set of $\bord^2\G$.
\end{defi}

\begin{rema}\label{r.conicalgeodesic}
It is easy to verify that a sequence $\{\g_n\}_{n\in\N}$ converges conically to $x\in\bord\G$ if and only if it lies in an uniform neighborhood of any geodesic ray $(\alpha_n)_0^\infty$ converging to $x$, namely there exists $K>0$ and a subsequence $\{\alpha_{n_k}\}$ such that for all $k$ one has $d_\G(\alpha_{n_k},\g_{k})<K$.
\end{rema}



Given $\g\in\G$ we denote by $\cone(\g)$ the \emph{cone type} of $\g\in\G$, namely
\[
\cone(\g):=\left\{h\in\G\,\middle|\, d(e,\g h)=d(e,\g)+d(e,h)\right\}.
\]
Cannon~\cite{CannonCones} showed the set of cone types of a Gromov hyperbolic group is finite, see for example Bridson--Haefliger's book~\cite[p.~455]{BH}. We denote by $\cone_\infty(\g)\subset{\bord\G}$ the set of points $x$ that can be represented by a geodesic ray contained in $\cone(\g)$.

\setcounter{section}{7}\setcounter{prop}{0}
% Original author-published.tex lines 2609--2646

Either $\G$ has boundary homeomorphic to a circle, or it is a Kleinian group. In the first case we let
\[
\rho,\ovrho:\G\to\Diff^{1+\nu}(\circle)
\]
be H\"older conjugated to action of $\G$ on its boundary; if instead $\G<\PSL(2,\C)$ is a~Kleinian group we let $\rho,\ovrho:\G\to\PSL(2,\C)$ be two convex co-compact representations that lie in the same connected component of the subset of the character variety $\fkX(\G,\PSL(2,\C))$ consisting of convex co-compact representations.

We let $X$ be either the circle or $\bord\H^3$. To simplify notation we will denote the action of $\g\in\G$ on $X$ via $\rho$ by $\g$, the action via $\ovrho$ by $\ov{\!\g}$, and the limit sets of $\rho$ and $\ovrho$ by $\bord\G,\bar{\bord\G}\subset X$ respectively.

In both situations there exists a H\"older-continuous map
\[
\Xi:X\to X
\]
conjugating $\rho$ and $\ovrho$. Indeed while in the surface case this holds by definition, in the Kleinian case this is a theorem by Marden~\cite{Marden}, see also Anderson's survey~\cite[p.~32]{SurveyAnderson}: the equivariant limit map $\Xi:\bord\G\to\bord\G$ conjugating the actions $\rho$ and $\ovrho$ on their respective limit sets extends to a $\G$-equivariant, H\"older continuous homeomorphism of the whole Riemann sphere $\bord\H^3$. We study differentiability points of $\Xi$ with oblique derivative.

We let $d$ be either a visual distance on $X$ (in the complex case) or a distance inducing the chosen $\class^1$ structure on the circle $\circle$.

\begin{defi}\label{assuA}
An action $\rho$ admits a \emph{Lipschitz-compatible cover} if there exists a finite open cover $\cal B$ of $X$ and a map $\G\to\cal B$, $\gamma\mapsto\cal B_\infty(\g)$ such that:
\begin{enumerate}\romanenumi
\item\label{defi7.1.1} for any $a,b\in \G$ so that $|ab|=|a|+|b|$ one has
\begin{enumerate}
\item\label{defi7.1.1.a} $b\cal B_{\infty}(ab)\subset \cal B_\infty(a)$,
\item\label{defi7.1.1.b} $\cal B_{\infty}(ab)\subset\cal B_{\infty}(b)$;
\end{enumerate}
\item\label{defi7.1.2} there exist $\lambda>0$, $C$ and $L\in\N$ such that if
$|\g|\geq L$ and $x,y\in\cal B_\infty(\g)$ then
\[
d(\g x,\g y)\leq Ce^{-|\g|\lambda}d(x,y);
\]
\item\label{defi7.1.3} there exist constants $r_1,r_2$ and a function $\slroot:\G\to\RR$ with $\slroot(\g)\geq\lambda|\g|$ such that for every $\g\in\G$ and every $x\in\cone_\infty(\g)$,
\[
B\left(x,r_1e^{-\slroot(\g)}\right)\subset\g\cal B_\infty(\g)\subset B\left(x, r_2e^{-\slroot(\g)}\right).
\]
\end{enumerate}
\end{defi}

The goal of the subsection is to prove the following result, similar arguments can be found in Guizhen~\cite{cui} in the context of conjugacies of expanding circle maps.

% Original author-published.tex lines 2649--2651
\begin{prop}\label{c.indepPer}
Let $\rho,\ovrho$ be as above and assume both admit a Lipschitz compatible cover. If there exists $p\in\bord\G$ such that $\Xi$ has a finite non-vanishing derivative (complex derivative in the Kleinian case) at $p$ then $\Xi|\bord\G$ is bi-Lipschitz.
\end{prop}

% Original author-published.tex lines 2655--2660
We work under the assumptions of Proposition~\ref{c.indepPer} and begin its proof with the following lemma. For $\g\in\G$ we denote its derivative at $x\in X$ by $\g'(x)\in\K$ defined, according our two situations, by:
\begin{itemize}
\item $X=\bbS^1$: the derivative $\tilde\g'(\tilde x)$ of a lift of $\g$ to the universal cover $\R$ of $\bbS^1$, and a lift $\tilde x\in\R$ of $x$, the number $\tilde \g'(\tilde x)$ is independent of these choices;
\item $X=\bord\H^3$: we fix an arbitrary point $\infty\notin\bord\G$, identify $X-\{\infty\}$ with $\C$ via the stereographic projection and let $\g'(x)$ be the standard complex derivative.
\end{itemize}


% Original author-published.tex lines 2662--2722
\begin{lemm}\label{controlC}
Let $\rho:\G\to\Diff^{1+\nu}(X)$ admit a Lipschitz compatible cover. There exists a constant $\kappa>0$ and $N\in\N$ such that for all $\g\in\G$ with $|\g|\geq N$ and $x,y\in\cal B_\infty(\g)$ one has
\[
\bigl\lvert\log\left|\g'(x)\right|-\log\left|\g'(y)\right|\bigr|\leq \kappa d(x,y)^\nu.
\]
\end{lemm}



\begin{proof}
We consider $L$ from Definition~\ref{assuA}, so that for every $\eta\in\G$ with $|\eta|\geq L$ and $x,y\in\cal B_\infty(\eta)$ one has
\begin{equation}\label{lip}
d(\eta x,\eta y)\leq Ce^{-|\eta|\lambda}d(x,y).
\end{equation}

Since the action is $\class^{1+\nu}$ we can find a positive $K$ such that for every $\beta$ with $|\beta|\leq L$ and $u,w\in X$ one has
\begin{equation}\label{derivativeGen}
\bigl\lvert\log\left|\beta'(u)\right|-\log\left|\beta'(w)\right|\bigr|\leq Kd(u,w)^\nu.
\end{equation}

We let then $K'=\max\{K,KC^\nu\}$. We begin by showing, by induction on $k$, that if $|\g|=kL$ then for all $x,y\in\cal B_\infty(\g)$, one has
\begin{equation}\label{induc}
\bigl\lvert\log\left|\g'(x)\right|-\log\left|\g'(y)\right|\bigr|\leq K'\left(\sum_{i=0}^{k-1}e^{-\nu\lambda L i}\right)d(x,y)^\nu.
\end{equation}

Eq.~\eqref{derivativeGen} gives the base case, so assume that the result holds up to $k-1$. We write $\g=\beta\eta$ with $|\beta|=L$, $|\eta|=(k-1)L$. By Definition~\ref{assuA} (ib) we have
\begin{equation}\label{conocont}
\cal B_\infty(\g)\subseteq\cal B_\infty(\eta).
\end{equation}
Applying the chain rule gives that for every $u\in X$ one has
\[
\log\left|\g'(u)\right|=\log\left|(\beta)'(\eta u)\right|+\log\left|(\eta)'(u)\right|
\]
and thus, when $x,y\in\cal B_\infty(\g)$,
\begin{multline*}
\bigl\lvert\log\left|\g'(x)\right|-\log\left|\g'(y)\right|\bigr\rvert\\
\begin{aligned}
& \leq \bigl\lvert\log\left|\beta'(\eta x)\right|-\log\left|\beta'(\eta y)\right|\bigr\rvert+\bigl|\log\left|\eta'(x)\right| -\log\left|\eta'(y)\right|\bigr|\\
& \leq K d(\eta x,\eta y)^\nu + K'\left(\sum_{i=0}^{k-2}e^{-\nu\lambda L i}\right)d(x,y)^\nu\quad (\text{by~\eqref{derivativeGen} and induction})\\
& \leq KC^\nu e^{-|\eta|\nu\lambda}d(x,y)^\nu+ K'\left(\sum_{i=0}^{k-2}e^{-\nu\lambda L i}\right)d(x,y)^\nu\quad (\text{by~\eqref{conocont} and~\eqref{lip}}).
\end{aligned}
\end{multline*}

This shows Eq.~\eqref{induc} which implies that for $\kappa_0=K'/(1-e^{-\nu\lambda L})$, every $\g\in\G$ whose word-length is an integer multiple of $L$, and $x,y\in\cal B_\infty(\g)$ one has
\[
\bigl\lvert\log\left|\g'(x)\right|-\log\left|\g'(y)\right|\bigr|\leq \kappa_0 d(x,y)^\nu.
\]


To conclude the lemma we consider an arbitrary $\g$ with $|\g|=mL+t$ and $t< L$. We write $\g=\beta \eta$ with $|\beta|=mL$. By Definition~\ref{assuA}$\MK$\eqref{defi7.1.1.a} it holds
\begin{equation}\label{e.6.6}
\eta \cal B_\infty(\g)\subset\cal B_\infty(\beta).
\end{equation}
Applying the chain rule gives then
\begin{align*}
\bigl\lvert\log\left|\g'(x)\right|-\log\left|\g'(y)\right|\bigr|& \leq \bigl\lvert\log\left|\beta'(\eta x)\right|-\log\left|\beta'(\eta y)\right|\bigr| +\bigl\lvert\log\left|\eta'(x)\right| -\log\left|\eta'(y)\right|\bigr|\\
& \leq \kappa_0d(\eta x,\eta y)^\nu + Kd(x,y)^\nu \quad (\text{by~\eqref{derivativeGen} and~\eqref{e.6.6}})\\
& \leq \left(\kappa_0C^\nu e^{-mL\lambda}+K\right) d(x,y)^\nu \;\:(\text{by~\eqref{lip}})
\end{align*}
so taking $\kappa=K+\kappa_0C^\nu e^{-L\lambda}$ we conclude the proof of Lemma~\ref{controlC}.
\end{proof}

% Original author-published.tex lines 2726--2789
\begin{proof}[Proof of Proposition~\ref{c.indepPer}]
Let $p\in\bord\G$ be such that $\Xi$ has a finite non-vanishing derivative. Fix a geodesic ray $(\alpha_n)_{0}^\infty$ through the identity with $\alpha_{n}\to p$. By definition for all $n$ one has $p\in\alpha_n\cone_\infty(\alpha_n)$. Without loss of generality we may also assume that
\[
p=0=\Xi(0)
\]
and we may write the derivative as the incremental limit
\[
\Xi'(0)=\lim_{y\to 0}\frac{\Xi(y)}{y}\in\K-\{0\}.
\]
For each $n$ we let $s_n=r_1e^{-\slroot(\alpha_n)}$, so that by Definition~\ref{assuA}$\MK$\eqref{defi7.1.3},
\[
B(0,s_n)\subset\alpha_n\cal B_\infty (\alpha_n).
\]

\pagebreak{}
We consider the scaling map
\[
g_n:B(0,1)\to \alpha_n\cal B_\infty(\alpha_n)
\]
defined by $g_n(z)=s_n z$.

Let $a_n$ be an arbitrary point at distance $s_n$ from $0$ and let $\tilde s_n=\Xi(a_n)$. Observe that since $\Xi$ is differentiable at zero, for $n$ big enough the image $\Xi(B(0,s_n))$ is coarsely a ball around zero of size comparable to that of $\ov\alpha_n\cone_\infty(\ov\alpha_n)$, and in particular we can assume, since the cover $\{\cal B_\infty(\ov{\!\g})\}$ is Lipschitz compatible (Definition~\ref{assuA}$\MK$\eqref{defi7.1.3}, that $\Xi(B(0,s_n))$ is contained in $\ov\alpha_n\cal B_\infty(\ov\alpha_n)$. Furthermore we deduce that there exist positive constants $d,D$ such that for every $n$
\[
d<\frac{\ov r_2e^{-\ovsloot(\ov\alpha_n)}}{|\tilde s_n|}<D.
\]

Here we denote by $\ov r_i, \ov \lambda, \bar{C}, \ovsloot$ the constants and function associated to the Lipschitz compatible cover $\{\cal B_\infty(\ov{\!\g})\}$ for the action $\ov \rho$. We consider the scaling map
\[
\tilde g_n:B(0,D)\to B(0,|\tilde s_n|D)
\]
by $z\mapsto z\tilde s_n$.

Since $s_n\to0$ and $\Xi'(0)\notin\{0,\infty\}$ exists, the composition
\[
\tilde g_n^{-1}\Xi g_n(z)=\frac{\Xi(zs_n)}{\tilde s_n}\cdot\frac{s_nz}{s_nz}=\frac{\Xi(zs_n)}{s_nz}\cdot\frac{s_n}{\tilde s_n}\cdot z=\frac{\Xi(zs_n)}{s_nz}\cdot\frac{s_n}{\Xi(s_n)}\cdot z
\]
converges uniformly on compact subsets to the identity map.



On the other hand, one has
\[
\tilde g_n^{-1}\Xi g_n=\tilde g_n^{-1}\overline\alpha_{n}\Xi\alpha_{n}^{-1} g_n.
\]
We now study the maps $f_n:=\alpha_{n}^{-1}\circ g_n$ and $\tilde f_n:=\tilde g_n^{-1}\circ\overline\alpha_{n}$. Since the coverings $\cal B$ and $\ov {\cal B}$ are finite, we can assume, up to extracting a subsequence that there exists sets $\cal B_\infty\in \cal B$, $\ov{\cal B}_\infty\in \ov {\cal B}$ so that, for every $n$, $\cal B_\infty(\alpha_n)=\cal B_\infty$ (resp. $\cal B_\infty(\ov\alpha_n)=\ov{\cal B}_\infty$).


Observe that for every $x\in B(0,1)$ one has
\[
\log\left|f_n'(x)\right|=\log\left|\left(\alpha_n^{-1}\right)'(g_nx)\right|+\log|s_n|= -\log\left|\alpha_n'\left(\alpha_n^{-1}g_nx\right)\right|+\log|s_n|.
\]
Now by definition of $g_n$, we have that $g_nx \in\alpha_n\cal B_\infty(\alpha_n)$ and thus $\alpha_n^{-1}(g_nx)\in\cal B_\infty(\alpha_n)$. For $n$ large enough we can apply Lemma~\ref{controlC} to $\alpha_n$ to obtain $\kappa$ so that for every pair $x,y\in B(0,1)$ it holds
\[
\bigl\lvert\log\left|f_n'(x)\right|-\log\left|f_n'(y)\right|\bigr|\leq \kappa d(x,y)^\nu.
\]
We conclude that the family of maps $(f_n)$ is uniformly bi-Lipschitz on $B(0,1)$ and thus, since $(f_n0)$ is bounded, Arzela--Ascoli's theorem applies to give a subsequence (still denoted by $f_n$) that converges to a bi-Lipschitz map $f$ defined on $B(0,1)$.


A similar reasoning applies to the maps $\tilde f$ defined on $\ov{\cal B}_\infty$, and we obtain that, about~$0$, $\Xi$ can be written as a composition of bi-Lipschitz maps and is thus bi-Lipschitz. Let $U$ be the open neighborhood where $\Xi$ is bi-Lipschitz, then for every $\g\in\G$, $\Xi$ is bi-Lipschitz on $\g U$ with possibly different constants. However, by compactness we can cover $\bord\G$ by finitely many open sets such that $\Xi$ is bi-Lipschitz restricted to each one of them. A priori the bi-Lipschitz constants of $\Xi$ depend on each open set but we may choose them to work simultaneously for these finitely many translates of~$U$. It remains then to check the bi-Lipschitz property on points which are uniformly far apart, i.e. whose distance is $\geq\eps$ for some $\eps>0$. However, the function
\[
(x,y)\mapsto \frac{d(f(x),f(y))}{d(x,y)}
\]
is continuous and non-vanishing on $\{(x,y):d(x,y)\geq\eps\}$. Since the latter space is compact we obtain the bi-Lipschitz property for points which are far apart, as desired.
\end{proof}

